\documentclass[a4paper,11pt]{article}
\usepackage{exam-style-842-unit}
\examinehead{842 信息化概论 · 第07单元练习册}{练习二\quad 人工智能应用与发展}
\begin{document}

\unitcover{842 信息化概论 · 单元练习册}{第07单元\quad 人工智能技术}{练习二\quad 人工智能应用与发展（共 100 分）}{约 70 分钟}{人工智能的 15 类应用（难题求解、自动规划调度配置、机器定理证明、自动程序设计、机器翻译、智能控制、智能管理、智能决策、智能通信、智能仿真、智能人机接口、模式识别、数据挖掘与知识发现、机器博弈、智能机器人）；数据挖掘与知识发现；智慧农业典型场景；人工智能伦理与挑战；人工智能的发展方向与趋势。}

\begin{knowmap}
\textbf{应用总览（教材 9.6，共 15 类）：}难题求解 · 自动规划调度与配置 · 机器定理证明 · 自动程序设计 · 机器翻译 · 智能控制 · 智能管理 · 智能决策 · 智能通信 · 智能仿真 · 智能人机接口 · 模式识别 · 数据挖掘与知识发现 · 机器博弈 · 智能机器人\\[2pt]
\textbf{分类视角：}符号智能类（定理证明、专家系统、机器博弈等，重符号推理）与计算智能类（模式识别、数据挖掘、智能机器人等，重数值计算与算法）\\[2pt]
\textbf{关键关系：}数据挖掘（DM）与数据库中的知识发现（KDD）本质含义一样，只是侧重领域不同——DM 偏统计、数据分析、数据库和信息系统，KDD 偏人工智能和机器学习\\[2pt]
\textbf{机器博弈演进：}1956 年塞缪尔跳棋程序 $\rightarrow$ 1997 年 IBM“深蓝”击败国际象棋冠军（主要靠运算能力）$\rightarrow$ 2016 年 AlphaGo 战胜李世石（主要原理是深度学习）\\[2pt]
\textbf{智慧农业典型场景（教材第 9 章未展开，按大纲、真题与第 10 章补充）：}病虫害识别 · 产量预测 · 智能农机作业 · 精准灌溉；教材第 10 章 10.2.3（7）农业自动化可作支撑\\[2pt]
\textbf{挑战与趋势：}伦理与挑战（隐私、算法偏见、就业替代、安全可控、责任归属）；近期可能方向 模糊处理 · 并行化 · 神经网络 · 机器情感；发展趋势 功能仿真 · 机器参与 · 辅助人脑思维 · 帮助人学习
\end{knowmap}

\begin{keyterms}
\textbf{难题求解}：求解没有算法解，或有算法解但在现有机器上无法实施或无法完成的困难问题，如 NP 问题、旅行商问题、路径规划、调度、预测与诊断（教材 9.6.1）。\\
\textbf{自动规划、调度与配置}：自动规划是设计制定行动序列，调度是任务分配或安排（数学本质是给出两个集合间的映射），配置是设计合理的部件组合结构；三者均属约束满足问题（教材 9.6.2）。\\
\textbf{机器定理证明}：用计算机证明数学定理，是最典型的逻辑推理问题，方法有自然演绎法、判定法、定理证明器、计算机辅助证明（教材 9.6.3）。\\
\textbf{自动程序设计}：让计算机设计程序，给出程序要求的高级描述后自动生成目标程序，并进行程序的自动验证（教材 9.6.4）。\\
\textbf{机器翻译}：完全用计算机作为两种语言之间的翻译（教材 9.6.5）。\\
\textbf{智能控制}：把人工智能技术引入控制领域，建立智能控制系统，其核心在于高层控制即组织级控制（教材 9.6.6）。\\
\textbf{智能决策支持系统}：决策支持系统与人工智能（特别是专家系统）相结合的产物，既能定量分析又能定性分析，可处理半结构化和非结构化问题（教材 9.6.8）。\\
\textbf{模式识别}：利用计算机进行物体识别，物体一般指文字、符号、图形、图像、语音、声音及传感器信息等实体对象，不包括概念、思想、意识等抽象对象（教材 9.6.12）。\\
\textbf{数据挖掘（DM）}：从大量数据中提取隐含的、有价值知识的技术，主要用于统计、数据分析、数据库和信息系统等领域（教材 9.6.13）。\\
\textbf{数据库中的知识发现（KDD）}：与数据挖掘本质含义相同，但主要用于人工智能和机器学习等领域（教材 9.6.13）。\\
\textbf{机器博弈}：用计算机进行棋类等对弈，是人工智能最早的研究领域之一，代表成果有塞缪尔跳棋程序、“深蓝”和 AlphaGo（教材 9.6.14）。\\
\textbf{智能机器人}：具备感觉、运动、思考三个要素，能感知环境并作出反应性动作的机器人，是人工智能技术的综合应用（教材 9.6.15）。\\
\textbf{智慧农业}：以信息和知识为核心要素，综合运用物联网、大数据、人工智能等技术实现农业生产、经营、管理、服务智能化的现代农业形态（大纲与真题口径，教材第 9 章未展开）。\\
\textbf{精准农业}：根据田间肥力、墒情、苗情、杂草、病虫害等信息，自动进行精确或精准的耕作、播种、施肥、灌溉、除草、喷药、收割等作业的农业生产方式（教材 10.2.3，印刷页 179—180）。
\end{keyterms}

\begin{contrast}{数据挖掘 / 数据库中的知识发现（每行依次为：核心含义·主要领域·相互关系）}
\citem{数据挖掘（DM）}{从大量数据中提取隐含的、有价值的知识}{统计、数据分析、数据库和信息系统等领域}{与 KDD 本质含义一样，可同时列出}
\citem{数据库中的知识发现（KDD）}{从数据库中提取隐含的、事先未知的、有潜在价值的知识}{人工智能和机器学习等领域}{与 DM 本质含义一样，侧重领域不同}
\end{contrast}

\begin{contrast}{传统决策支持系统 / 智能决策支持系统（每行依次为：分析能力·处理问题类型·构成方式）}
\citem{传统决策支持系统}{以数值分析为主，擅长定量分析}{主要处理结构化问题}{模型库、方法库、数据库等}
\citem{智能决策支持系统}{定量分析与定性分析相结合}{可处理半结构化和非结构化问题}{传统决策支持系统加上智能部件，如专家系统模式、知识库系统模式}
\end{contrast}

\begin{contrast}{“深蓝” / 阿尔法围棋（AlphaGo）（每行依次为：主要原理·技术特征·代表意义）}
\citem{“深蓝”}{主要靠运算能力，输入数百万个国际象棋案例后筛选最佳着法}{以高速搜索和案例匹配为主}{1997 年击败国际象棋世界冠军卡斯帕罗夫}
\citem{阿尔法围棋（AlphaGo）}{主要工作原理是深度学习}{多层神经网络、海量对局训练、自学与自我对弈}{2016 年以 4:1 战胜世界围棋冠军李世石}
\end{contrast}

\begin{contrast}{智慧农业 / 精准农业（每行依次为：核心内涵·技术支撑·侧重点）}
\citem{智慧农业}{以信息和知识为核心要素，实现农业全链条智能化}{物联网、大数据、人工智能、云计算等}{覆盖生产、经营、管理、服务全过程}
\citem{精准农业}{按需精准作业，减少投入、提高效率与保护生态}{GPS、遥感、传感器、智能农机、自动控制等}{侧重田间作业的精准化与自动化}
\end{contrast}

\qtypename{一、单项选择题}{10 题，每题 3 分，共 30 分}
\q 教材把人工智能的应用归纳为（\quad）
\choicesline{A. 8 个方面}{B. 10 个方面}{C. 15 个方面}{D. 20 个方面}

\q 自动规划、调度与配置属于人工智能经典问题中的（\quad）
\choicesline{A. 约束满足问题}{B. 分类问题}{C. 聚类问题}{D. 降维问题}

\q 我国数学家吴文俊 1977 年提出的初等几何定理证明方法属于（\quad）
\choicesline{A. 自然演绎法}{B. 判定法}{C. 定理证明器}{D. 计算机辅助证明}

\q 智能控制的核心在于（\quad）
\choices{A. 低层执行机构的精度}{B. 高层控制，即组织级控制}
{C. 通信带宽的提升}{D. 传感器数量的增加}

\q 关于智能决策支持系统，下列说法正确的是（\quad）
\choices{A. 只能进行数值分析、定量分析}{B. 只能进行定性分析}
{C. 既能定量分析又能定性分析，可处理半结构化和非结构化问题}{D. 与专家系统无关，不需要知识}

\q 模式识别中所说的“物体”一般指（\quad）
\choices{A. 概念、思想、意识等抽象对象}{B. 文字、符号、图形、图像、语音、声音及传感器信息等实体对象}
{C. 仅指三维立体物体}{D. 仅指生物个体}

\q 关于数据挖掘（DM）与数据库中的知识发现（KDD）的关系，正确的是（\quad）
\choices{A. 二者本质含义一样，只是侧重领域不同}{B. 二者含义完全不同，毫无联系}
{C. KDD 只用于统计分析和数据库领域}{D. DM 只用于人工智能和机器学习领域}

\q 1997 年以 2 胜 3 平 1 负击败国际象棋世界冠军卡斯帕罗夫的 IBM 计算机是（\quad）
\choicesline{A. 沃森}{B.“深蓝”}{C. AlphaGo}{D. ENIAC}

\q 2016 年战胜围棋世界冠军李世石的 AlphaGo，其主要工作原理是（\quad）
\choicesline{A. 专家系统}{B. 深度学习}{C. 遗传算法}{D. 蚁群算法}

\q 按教材 9.6.15，智能机器人一般至少具备的三个要素是（\quad）
\choices{A. 感觉要素、运动要素、思考要素}{B. 感觉要素、反应要素、记忆要素}
{C. 输入要素、输出要素、存储要素}{D. 电源要素、驱动要素、控制要素}

\qtypename{二、多项选择题}{3 题，每题 4 分，共 12 分；每题有两个或两个以上正确项}
\q 下列各项中，属于人工智能典型应用领域的有（\quad）
\choicesline{A. 机器翻译}{B. 模式识别}{C. 机器博弈}{D. 人工合成淀粉}

\q 关于“深蓝”与 AlphaGo，下列说法正确的有（\quad）
\choices{A.“深蓝”的超人表现几乎纯粹靠运算能力}{B. AlphaGo 的主要工作原理是深度学习}
{C. AlphaGo 的核心是穷举围棋的所有落子可能}{D. 两者都是机器博弈的典型代表}

\q 关于智能机器人，下列说法正确的有（\quad）
\choices{A. 智能机器人一般至少具备感觉、运动、思考三个要素}{B. 智能机器人是人工智能技术的综合应用}
{C. 工业机器人已在结构化环境中执行各类确定性任务}{D. 智能机器人已经完全具备人类的意识和直觉}

\qtypename{三、判断题}{6 题，每题 2 分，共 12 分；正确的画“√”，错误的画“×”并写出改正后的表述}
\q 自动规划、调度与配置都属于人工智能经典问题中的约束满足问题。\quad（\quad）\quad 改正：\fixblank

\q 吴文俊提出的初等几何定理证明方法属于自然演绎法。\quad（\quad）\quad 改正：\fixblank

\q 智能决策支持系统只能进行定量分析，不能处理半结构化和非结构化问题。\quad（\quad）\quad 改正：\fixblank

\q 模式识别中的“物体”包括概念、思想等抽象对象。\quad（\quad）\quad 改正：\fixblank

\q 数据挖掘与数据库中的知识发现本质含义一样，只是侧重领域不同。\quad（\quad）\quad 改正：\fixblank

\q 智能机器人是人工智能技术的综合应用，其能力和水平是人工智能技术水平的一个体现。\quad（\quad）\quad 改正：\fixblank

\qtypename{四、填空题}{6 空，每空 1 分，共 6 分}
\q 教材 9.6 把人工智能的应用归纳为\fillblanks[1.2cm]个方面。

\q 1997 年，IBM 研制的计算机\fillblanks[2.0cm]击败了国际象棋世界冠军卡斯帕罗夫。

\q 2016 年 AlphaGo 战胜李世石，其主要工作原理是\fillblanks[2.2cm]。

\q 智能决策支持系统是决策支持系统与\fillblanks[2.2cm]（特别是专家系统）相结合的产物。

\q 数据挖掘与数据库中的知识发现的英文缩写分别是 DM 和\fillblanks[1.4cm]。

\q 教材 9.7.3 指出，未来人工智能可能会向模糊处理、并行化、神经网络和\fillblanks[2.2cm]四个方面发展。

\qtypename{五、简答题}{4 题，每题 6 分，共 24 分}
\q 简述人工智能的 15 个应用方面（至少列举 10 个，并各作一句话说明）。\anslines{2.4cm}

\q 简述智能决策支持系统的含义、智能部件的模式及其与传统决策支持系统的区别。\anslines{2.4cm}

\q 简述数据挖掘与知识发现的联系与区别，并说明当前人工智能技术形成的 5 条主线。\anslines{2.4cm}

\q 简述人工智能发展的方向与趋势（近期可能方向与发展趋势各写 4 点）。\anslines{2.4cm}

\newpage
\qtypename{六、综合题（论述）}{1 题，共 16 分}
\q 人工智能技术在农业上的应用是历年真题的高频考点。请结合教材第 9 章与考试大纲论述：（1）人工智能技术的主要内容（八大类）在农业中有哪些对应体现；（2）从人工智能的 15 类应用中至少举出 5 类，说明它们在农业领域的具体应用；（3）智慧农业的典型应用场景（至少 4 个）及其支撑技术；（4）人工智能在农业应用中面临的伦理与挑战；（5）人工智能的发展趋势及其对农业信息化的意义。\anslines{5cm}

\end{document}
